\section{符号与约定}

\subsection{向量的定义}

\begin{BoxDefinition}[向量的定义]
    $\vb{x}\in\R^n$：$\vb{x}$是一个$n$维的实向量（Real-Valued Vector）。
    \begin{Equation}[向量的定义]
        \vb{x}=\begin{bmatrix}
            x_1\\
            x_2\\
            \vdots\\
            x_n
        \end{bmatrix},\quad x_i\in\R
    \end{Equation}
    $\vb{x}\in\C^n$：$\vb{x}$是一个$n$维的复向量（Complex-Valued Vector）。
\end{BoxDefinition}

在线性代数的语境中，向量$\vb{x}$默认指的都是列向量（Column Vector），通过转置的记法$\vb{x}^T$表示行向量（Row Vector）。行内公式中，常用$\vb{x}=\qty[x_1, x_2, \cdots, x_n]^T$表示列向量以节约空间。

\subsection{矩阵的定义}

\begin{BoxDefinition}[实矩阵的定义]
    $\vb{A}\in\R^{m\times n}$：$\vb{A}$是一个$m\times n$维的实矩阵（Real-Valued Matrix）。
    \begin{Equation}[实矩阵的定义]
        \vb{A}=\begin{bmatrix}
            a_{11}&a_{12}&\cdots&a_{1n}\\
            a_{21}&a_{22}&\cdots&a_{2n}\\
            \vdots&\vdots&\ddots&\vdots\\
            a_{m1}&a_{m2}&\cdots&a_{mn}\\
        \end{bmatrix},\quad a_{ij}\in\R
    \end{Equation}
    $\vb{A}\in\C^{m\times n}$：$\vb{A}$是一个$m\times n$维的复矩阵（Complex-Valued Matrix）。
\end{BoxDefinition}

矩阵$\vb{A}\in\R^{m\times n}$称为$m$行$n$列，元素$a_{ij}$的下标代表其是$\vb{A}$的第$i$行第$j$列处的元素。

矩阵可以视为若干列向量横向拼接的结果，例如$\vb{A}\in\R^{m\times n}$可以记为
\begin{Equation}[矩阵的列向量拼接]
    \vb{A}=\qty[\vb{a_1},\vb{a_2},\cdots,\vb{a_n}]
\end{Equation}

其中$\vb{a}_j\in\R^m$代表代表矩阵$\vb{A}$的第$j$列。

\subsection{转置}

\begin{BoxDefinition}[转置]
    对于矩阵$\vb{A}\in\R^{m\times n}$，定义其转置$\vb{A^T}\in\R^{n\times m}$（Transpose）
    \begin{Equation}[转置]
        \vb{A}^T=\begin{bmatrix}
            a_{11}&a_{21}&\cdots&a_{n1}\\
            a_{12}&a_{22}&\cdots&a_{n2}\\
            \vdots&\vdots&\ddots&\vdots\\
            a_{1m}&a_{2m}&\cdots&a_{nm}\\
        \end{bmatrix}
    \end{Equation}
\end{BoxDefinition}

转置的实质就是将矩阵的行和列对调，原先的行变成了列，原先的列变成了行。

\begin{BoxProperty}[矩阵数乘的转置]
    矩阵数乘的转置，等于矩阵转置的数乘
    \begin{Equation}[矩阵数乘的转置]
        \qty(c\vb{A})^T=c\vb{A}^T
    \end{Equation}
\end{BoxProperty}

\begin{BoxProperty}[矩阵和的转置]
    矩阵和的转置，等于矩阵转置的和
    \begin{Equation}[矩阵和的转置]
        \qty(\vb{A}+\vb{B})^T=\vb{A}^T+\vb{B}^T
    \end{Equation}
\end{BoxProperty}

\begin{BoxProperty}[矩阵积的转置]
    矩阵积的转置，等于矩阵转置的逆序积
    \begin{Equation}[矩阵积的转置]
        \qty(\vb{A}\vb{B})^T=\vb{B}^T\vb{A}^T
    \end{Equation}
\end{BoxProperty}

转置计算可以提到矩阵乘积内进行，但务必记得交换乘法顺序！（矩阵乘法不满足交换律）

\begin{BoxProperty}[双重转置还原]
    矩阵转置的转置等于原矩阵自身
    \begin{Equation}[双重转置还原]
        (\vb{A}^T)^T=\vb{A}
    \end{Equation}
\end{BoxProperty}

\begin{BoxDefinition}[对称矩阵]
    若$\vb{A}\in\R^{n\times}$满足以下性质，则称为对称矩阵（Symetric Matrix）
    \begin{Equation}[对称矩阵]
        \vb{A}^T=\vb{A}
    \end{Equation}
    若$\vb{A}\in\R^{n\times}$满足以下性质，则称为协对称矩阵（Skew-Symetric Matrix）
    \begin{Equation}[协对称矩阵]
        \vb{A}^T=-\vb{A}
    \end{Equation}
\end{BoxDefinition}

\begin{BoxProperty}[对称分解]
    对于任意$\vb{A}\in\R^{n\times n}$，其可以被分解为
    \begin{Equation}[对称分解]
        \vb{A}=\vb{T}+\vb{S}
    \end{Equation}
    其中，$\vb{T}=(\vb{A}+\vb{A}^T)/2$是对称矩阵，$\vb{S}=(\vb{A}-\vb{A}^T)/2$是协对称矩阵。
\end{BoxProperty}

\subsection{伴随}

\begin{BoxDefinition}[伴随]
    对于矩阵$\vb{A}\in\C^{m\times n}$，定义其伴随$\vb{A^H}\in\C^{n\times m}$（Adjoint）
    \begin{Equation}[伴随]
        \vb{A}^T=\begin{bmatrix}
            a_{11}^{*}&a_{21}^{*}&\cdots&a_{n1}^{*}\\
            a_{12}^{*}&a_{22}^{*}&\cdots&a_{n2}^{*}\\
            \vdots&\vdots&\ddots&\vdots\\
            a_{1m}^{*}&a_{2m}^{*}&\cdots&a_{nm}^{*}\\
        \end{bmatrix}
    \end{Equation}
\end{BoxDefinition}

伴随也可以称为共轭转置（Conjugate Transpose）或厄米转置（Hermitian Transpose）。

伴随的实质是：对复矩阵求转置后对每一元素求复共轭。

\begin{BoxProperty}[矩阵数乘的伴随]
    矩阵数乘的转置，等于矩阵伴随的数乘
    \begin{Equation}[矩阵数乘的伴随]
        \qty(c\vb{A})^H=c\vb{A}^H
    \end{Equation}
\end{BoxProperty}

\begin{BoxProperty}[矩阵和的伴随]
    矩阵和的伴随，等于矩阵伴随的和
    \begin{Equation}[矩阵和的伴随]
        \qty(\vb{A}+\vb{B})^H=\vb{A}^H+\vb{B}^H
    \end{Equation}
\end{BoxProperty}

\begin{BoxProperty}[矩阵积的伴随]
    矩阵积的伴随，等于矩阵伴随的逆序积
    \begin{Equation}[矩阵积的伴随]
        \qty(\vb{A}\vb{B})^H=\vb{B}^H\vb{A}^H
    \end{Equation}
\end{BoxProperty}

\begin{BoxProperty}[双重伴随还原]
    矩阵伴随的伴随等于原矩阵自身
    \begin{Equation}[双重伴随还原]
        (\vb{A}^H)^H=\vb{A}
    \end{Equation}
\end{BoxProperty}

\begin{BoxDefinition}[厄米矩阵]
    若$\vb{A}\in\R^{n\times}$满足以下性质，则称为厄米矩阵（Hermitian Matrix）
    \begin{Equation}[厄米矩阵]
        \vb{A}^T=\vb{A}
    \end{Equation}
    若$\vb{A}\in\R^{n\times}$满足以下性质，则称为协厄米矩阵（Skew-Hermitian Matrix）
    \begin{Equation}[协厄米矩阵]
        \vb{A}^T=-\vb{A}
    \end{Equation}
\end{BoxDefinition}

\begin{BoxProperty}[厄米分解]
    对于任意$\vb{A}\in\R^{n\times n}$，其可以被分解为
    \begin{Equation}[厄米分解]
        \vb{A}=\vb{T}+\vb{S}
    \end{Equation}
    其中，$\vb{T}=(\vb{A}+\vb{A}^T)/2$是厄米矩阵，$\vb{S}=(\vb{A}-\vb{A}^T)/2$是协厄米矩阵。
\end{BoxProperty}

\subsection{迹}

\begin{BoxDefinition}[迹]
    对于矩阵$\vb{A}\in\R^{n\times n}$，定义其迹（Trace）
    \begin{Equation}[迹]
        \tr(\vb{A})=\Sum[i=1][n]a_{ii}
    \end{Equation}
\end{BoxDefinition} 

简而言之，矩阵的迹就是矩阵主对角线上的元素之和！

\begin{BoxProperty}[矩阵转置的迹]
    矩阵及其转置具有相同的迹
    \begin{Equation}[矩阵转置的迹]
        \tr(\vb{A}^T)=\tr(\vb{A})
    \end{Equation}
\end{BoxProperty}

\begin{BoxProperty}[矩阵数乘的迹]
    矩阵数乘的迹，等于矩阵迹的数乘
    \begin{Equation}[矩阵数乘的迹]
        \tr(c\vb{A})=c\tr(\vb{A})
    \end{Equation}
\end{BoxProperty}

\begin{BoxProperty}[矩阵和的迹]
    矩阵和的迹，等于矩阵迹的和   
    \begin{Equation}[矩阵和的迹]
        \tr(\vb{A}+\vb{B})=\tr(\vb{A})+\tr(\vb{B})
    \end{Equation}
\end{BoxProperty}

\begin{BoxProperty}[矩阵积的迹]
    矩阵的相乘顺序，不影响乘积的迹
    \begin{Equation}[矩阵积的迹]
        \tr(\vb{A}\vb{B})=\tr(\vb{B}\vb{A})
    \end{Equation}
\end{BoxProperty}

\begin{BoxProperty}[对称与协对称矩阵的积的迹]
    若$\vb{A}$是对称矩阵，而$\vb{B}$是协对称矩阵，则其积的迹恒为零
    \begin{Equation}[对称与协对称矩阵的积的迹]
        \tr(\vb{A}\vb{B})=0
    \end{Equation}
\end{BoxProperty}

\subsection{幂}

\begin{BoxDefinition}[矩阵的幂]
    对于矩阵$\vb{A}\in\R^{n\times n}$，记号$\vb{A}^k$定义为
    \begin{Equation}[矩阵的幂]
        \vb{A}^k=\underbrace{\vb{A}\vb{A}\cdots\vb{A}}_{k\ \text{times}}
    \end{Equation}
\end{BoxDefinition}

\subsection{特殊的向量}

\begin{BoxDefinition}[零向量]
    记号$\vb{0}$表示一个所有元素都为$0$的零向量（Zero Vector）
    \begin{Equation}[零向量]
        \vb{0}=\begin{bmatrix}
            0\\
            \vdots\\
            0\\
        \end{bmatrix}
    \end{Equation}
\end{BoxDefinition}

需要指出的是，记号$\vb{0}$不仅可以用来表示零向量，也可以用来表示零矩阵。

\begin{BoxDefinition}[全一向量]
    记号$\vb{1}$表示一个所有元素都为$1$的全一向量（All-One Vector）
    \begin{Equation}[全一向量]
        \vb{1}=\begin{bmatrix}
            1\\
            \vdots\\
            1\\
        \end{bmatrix}
    \end{Equation}
\end{BoxDefinition}

\begin{BoxDefinition}[单位向量]
    记号$\vb{e}_i$表示仅有第$i$个元素为$1$的单位向量（Unit Vector）
    \begin{Equation}[单位向量]
        \vb{e}_i=\begin{bmatrix}
            0\\
            \vdots\\
            1\\
            \vdots\\
            0\\
        \end{bmatrix}
    \end{Equation}
\end{BoxDefinition}

\subsection{特殊的矩阵}

\begin{BoxDefinition}[单位矩阵]
    记号$\vb{I}$表示仅有主对角线元素为$1$的单位矩阵（Identity Matrix）
    \begin{Equation}[单位矩阵]
        \vb{I}=\begin{bmatrix}
            1 &        &  \\
              & \ddots &  \\
              &        & 1\\
        \end{bmatrix}
    \end{Equation}
\end{BoxDefinition}

\begin{BoxDefinition}[对角矩阵]
    记号$\diag(a_1,\cdots,a_n)$表示仅有主对角线元素的对角矩阵（Diagonal Matrix）
    \begin{Equation}[对角矩阵]
        \diag(a_1,\cdots,a_n)=\begin{bmatrix}
            a_1 &        &     \\
                & \ddots &     \\
                &        & a_n \\
        \end{bmatrix}
    \end{Equation}
\end{BoxDefinition}

若对角元素表示为向量$\vb{a}=[a_1,\cdots,a_n]^T$，也可以将$\diag(a_1,\cdots,a_n)$简记为$\diag(\vb{a})$。
\begin{itemize}
    \item 将满足$i\leq j$，位于主对角线上方的$a_{ij}$称为超对角线（Super-Diagonal）。
    \item 将满足$i\geq j$，位于主对角线下方的$a_{ij}$称为次对角线（Sub-Diagonal）。
\end{itemize}

\subsection{矩阵的尺寸}

\begin{BoxDefinition}[高矩阵]
    对于矩阵$\vb{A}\in\R^{m\times n}$，若$m>n$，称$\vb{A}$为高矩阵（Tall）。
\end{BoxDefinition}

\begin{BoxDefinition}[宽矩阵]
    对于矩阵$\vb{A}\in\R^{m\times n}$，若$m<n$，称$\vb{A}$为宽矩阵（Tall）。
\end{BoxDefinition}

简而言之，行比列多的矩阵是高矩阵，行比列少的矩阵是宽矩阵。

\begin{BoxDefinition}[方阵]
    对于矩阵$\vb{A}\in\R^{m\times n}$，若$m=n$，称$\vb{A}$为方阵（Square）。
\end{BoxDefinition}

\subsection{矩阵的三角性}

\begin{BoxDefinition}[上三角矩阵]
    对于矩阵$\vb{A}\in\R^{n\times n}$，若$i>j$皆有$a_{ij}=0$，称$\vb{A}$为上三角矩阵（Upper Traiangular）。
    \begin{Equation}[上三角矩阵的例子]
        \vb{A}=\begin{bmatrix}
            a_{11} & a_{12} & \cdots & a_{1n} \\
            0      & a_{22} & \cdots & a_{2n} \\
            \vdots & \vdots & \ddots & \vdots \\
            0      & 0      & \cdots & a_{nn} \\
        \end{bmatrix}
    \end{Equation}
    若放宽至$i>j+1$，允许次对角线非零，称$\vb{A}$为上海森堡矩阵（Upper Hessenberg）。
\end{BoxDefinition}

\begin{BoxDefinition}[下三角矩阵]
    对于矩阵$\vb{A}\in\R^{n\times n}$，若$i<j$皆有$a_{ij}=0$，称$\vb{A}$为下三角矩阵（Lower Traiangular）。
    \begin{Equation}[下三角矩阵的例子]
        \vb{A}=\begin{bmatrix}
            a_{11}&0     &\cdots&0     \\
            a_{21}&a_{22}&\cdots&0     \\
            \vdots&\vdots&\ddots&\vdots\\
            a_{n1}&a_{n2}&\cdots&a_{nn}\\
        \end{bmatrix}
    \end{Equation}
    若放宽至$i<j-1$，允许超对角线非零，称$\vb{A}$为上海森堡矩阵（Upper Hessenberg）。
\end{BoxDefinition}

上三角矩阵和下三角矩阵讨论的都是方阵，其“上”和“下”指的是有元素的部分。

\begin{BoxDefinition}[带状矩阵]
    对于矩阵$\vb{A}\in\R^{n\times n}$，若满足
    \begin{Equation}[带状矩阵]
        a_{ij}=0,\qquad i>j+p \qor j>i+q
    \end{Equation}
    则称为带状矩阵（Banded Matrix）。
\end{BoxDefinition}

带状矩阵是关于对角矩阵和三角矩阵的进一步统一，如\cref{tab:带状矩阵的特例}所示
\begin{Table}[带状矩阵的特例]!!
    \begin{tblr}{colspec={llllX[4,r]},cell{1-Z}{1,2}={mode=math}}
        p&q&名称&英文&哪些位置有元素？\\
        0&0&对角矩阵&Diagonal&主对角线\\
        1&1&三对角矩阵&Tridiagonal&主对角线、超对角线、次对角线\\
        0&1&上双对角矩阵&Upper Bidiagonal&主对角线、超对角线\\
        1&0&下双对角矩阵&Lower Bidiagonal&主对角线、次对角线\\
        0&n-1&上三角矩阵&Upper Traiangular&主对角线及以上\\
        n-1&0&下三角矩阵&Lower Traiangular&主对角线及以下\\
        1&n-1&上海森堡矩阵&Upper Hessenberg&次对角线及以上\\
        n-1&1&下海森堡矩阵&Lower Hessenberg&超对角线及以下\\
    \end{tblr}
\end{Table}

\subsection{结构特殊矩阵}

\begin{BoxDefinition}[托普利茨矩阵]
    定义托普利茨矩阵（Toeplitz Matrix）为具有恒定对角元素的矩阵
    \begin{Equation}[托普利茨矩阵]
        \vb{A}=\mqty[
            a_{0}   &a_{-1}     &a_{-2}     &\cdots     &\cdots     &a_{-(n-1)} \\
            a_{1}   &a_{0}      &a_{-1}     &\ddots     &           &\vdots     \\
            a_{2}   &a_{1}      &\ddots     &\ddots     &\ddots     &\vdots     \\
            \vdots  &\ddots     &\ddots     &\ddots     &a_{-1}     &a_{-2}     \\
            \vdots  &           &\ddots     &a_{1}      &a_{0}      &a_{-1}     \\
            a_{n-1} &\cdots     &\cdots     &a_{2}      &a_{1}      &a_{0}      \\
        ]
    \end{Equation}
    特别的，若$a_{-k}=a_{n-k}$，则称为循环矩阵（Circulant Matrix）。
\end{BoxDefinition}

若托普利茨矩阵上下颠倒，就是所谓的汉克尔矩阵（恒定元素由对角元素变为反对角元素）。

\begin{BoxDefinition}[汉克尔矩阵]
    定义汉克尔矩阵（Hankel Matrix）为具有恒定反对角元素的矩阵
    \begin{Equation}[汉克尔矩阵]
        \vb{A}=\mqty[
            a_{0}   &\cdots     &\cdots     &a_{n-3}    &a_{n-2}    &a_{n-1}    \\
            a_{1}   &           &\adots     &a_{n-2}    &a_{n-1}    &a_{n}      \\
            \vdots  &\adots     &\adots     &\adots     &a_{n}      &a_{n+1}    \\
            a_{n-3} &a_{n-2}    &\adots     &\adots     &\adots     &\vdots     \\
            a_{n-2} &a_{n-1}    &a_n        &\adots     &           &a_{2n-3}   \\
            a_{n-1} &a_{n}      &a_{n+1}    &\cdots     &\cdots     &a_{2n-2}   \\
        ]
    \end{Equation}
\end{BoxDefinition}

\begin{BoxDefinition}[范德蒙矩阵]
    定义范德蒙矩阵（Vandemonde Matrix）是满足$a_{ij}=\alpha_i^{j-1}$的矩阵
    \begin{Equation}[范德蒙矩阵]
        \vb{A}=\mqty[
            1       &\alpha_1   &\alpha_1^2     &\cdots     &\alpha_1^{n-1} \\
            1       &\alpha_2   &\alpha_2^2     &\cdots     &\alpha_2^{n-1} \\
            1       &\alpha_3   &\alpha_3^2     &\cdots     &\alpha_3^{n-1} \\
            \vdots  &\vdots     &\vdots         &\ddots     &\vdots         \\
            1       &\alpha_m   &\alpha_m^2     &\cdots     &\alpha_m^{n-1} \\
        ]
    \end{Equation}
\end{BoxDefinition}

\subsection{代数特殊矩阵}

\begin{BoxDefinition}[对合矩阵]
    定义对合矩阵（Involutory Matrix）满足以下条件
    \begin{Equation}[对合矩阵]
        \vb{A}^2=\vb{I}
    \end{Equation}
\end{BoxDefinition}

\begin{BoxDefinition}[幂等矩阵]
    定义幂等矩阵（Idempotent Matrix）满足以下条件
    \begin{Equation}[幂等矩阵]
        \vb{A}^2=\vb{A}
    \end{Equation}
\end{BoxDefinition}

\begin{BoxDefinition}[幂零矩阵]
    定义幂零矩阵（Nilpotent Matrix）为对于某个$k>0$，满足
    \begin{Equation}[幂零矩阵]
        \vb{A}^k=\vb{0}
    \end{Equation}
\end{BoxDefinition}

若一个矩阵$\vb{A}$满足$\vb{A}-\vb{I}$是幂零矩阵，那就称之为幺零矩阵。

\begin{BoxDefinition}[幺零矩阵]
    定义幺零矩阵（Unipotent Matrix）为对于某个$k>0$，满足
    \begin{Equation}[幺零矩阵]
        (\vb{A}-\vb{I})^k=\vb{0}
    \end{Equation}
\end{BoxDefinition}

\subsection{子矩阵}

\begin{BoxDefinition}[子矩阵]
    子矩阵（Submatrix）定义为从矩阵$\vb{A}$中选取由$\vb*{\alpha}$指定的行和由$\vb*{\beta}$指定的列。
\end{BoxDefinition}

\begin{BoxDefinition}[主子矩阵]
    主子矩阵（Principal Submatrix）的约束是选取的行和列要相同
    \begin{Equation}[主子矩阵]
        \vb*{\alpha}=\vb*{\beta}
    \end{Equation}
\end{BoxDefinition}

\begin{BoxDefinition}[前导主子矩阵]
    前导主子矩阵（Leading Principal Submatrix）的约束是选取从左上开始的连续行列
    \begin{Equation}[前导主子矩阵]
        \vb*{\alpha}=\vb*{\beta}=\qty[1,\cdots,k]^T
    \end{Equation}
\end{BoxDefinition}

\begin{BoxDefinition}[后导主子矩阵]
    后导主子矩阵（Trailing Principal Submatrix）的约束是选取从右下开始的连续行列
    \begin{Equation}[后导主子矩阵]
        \vb*{\alpha}=\vb*{\beta}=\qty[l,\cdots,\min\qty{m,n}]^T
    \end{Equation}
\end{BoxDefinition}
